\documentclass[11pt,a4paper]{article}

\usepackage[utf8]{inputenc}
\usepackage[T1]{fontenc}
\usepackage[english]{babel}
\usepackage{amsmath,amssymb,amsthm}
\usepackage{geometry}
\usepackage{microtype}
\usepackage{graphicx}
\usepackage{booktabs}
\usepackage{hyperref}
\usepackage{cite}

\newtheorem{theorem}{Theorem}

\geometry{margin=2.5cm}
\hypersetup{
  colorlinks=true,
  linkcolor=blue,
  citecolor=blue,
  urlcolor=blue
}

\title{Second Interim Report:\\[0.3em]
\large Real-Time Face Reenactment via SVD-Based Landmark Mapping}
\author{Viktor Syrotiuk \and Mykhailo Brytan \and Yulian Volodymyr Zaiats}
\date{\today}

\begin{document}
\maketitle

\begin{abstract}
This report extends the first interim submission on landmark-based face reenactment.
The objective is to evaluate whether a transparent linear-algebra pipeline can
transfer facial expressions from a live video stream to a static portrait in near
real time with acceptable geometric fidelity. Expression transfer is formulated as
a linear least-squares estimation problem solved via the singular value decomposition
(SVD). The document elaborates the problem setting, reviews related work, describes
the chosen algorithm with its advantages and limitations, develops the theoretical
background with emphasis on linear algebra methods, and presents the implementation
pipeline with completed and planned work. A link to the public GitHub repository is
included at the end.
\end{abstract}

\section{Introduction and problem statement}
\label{sec:intro}

\paragraph{General topic.}
The project lies at the intersection of \emph{computer vision} and \emph{numerical
linear algebra}. \emph{Face reenactment} refers to driving the appearance of one
face (e.g.\ a photograph) by the facial motion observed in another (e.g.\ webcam
input). Commercial and research systems often rely on deep generative models; the
present work instead adopts a \emph{geometry-first} formulation in which facial
geometry is summarized by 2D landmark points and the mapping between expressions is
estimated with a linear model and a robust linear solver.

\paragraph{Detailed problem.}
The setup involves:
\begin{itemize}
  \item a \textbf{source image}~$I_s$ (static portrait, e.g.\ a famous person);
  \item a \textbf{target video stream}~$\{I_t\}$ from a camera, capturing a
        performer whose expressions should be mimicked by the source face.
\end{itemize}
For each relevant frame a landmark detector yields $k$ facial key points. These are
stacked into a single vector
\begin{equation}
  x = (x_1,y_1,\ldots,x_k,y_k)^{\mathsf T} \in \mathbb{R}^{2k}.
  \label{eq:landmark-vec}
\end{equation}
Similarly, $y \in \mathbb{R}^{2k}$ denotes the target landmark vector. During
a~\emph{calibration} phase several pairs $(x^{(i)}, y^{(i)})$ are collected. The
unknown parameters~$a$ of a linear transformation model encode how to predict
target-aligned landmarks from the source configuration. Stacking the calibration
equations produces an overdetermined linear system
\begin{equation}
  M a \approx b,
  \label{eq:Ma-b}
\end{equation}
where $M \in \mathbb{R}^{m \times n}$ ($m > n$) and $b \in \mathbb{R}^{m}$ are
assembled from the landmark correspondences. The estimate $\hat{a}$ minimises
$\|Ma-b\|_2^2$. At runtime, $\hat{a}$ is applied together with freshly detected
target landmarks to compute predicted source landmarks; an image warp maps the
pixels of~$I_s$ to produce the reenacted frame.

\paragraph{Research aim.}
The aim of the project is to \textbf{implement and assess} how far a transparent,
landmark-based linear reenactment pipeline can go---in geometric fidelity and
latency---compared to a neural baseline, while keeping the core mathematics
(least squares, orthogonality, SVD) explicit and reproducible.

\section{Related work and possible approaches}
\label{sec:related}

Face reenactment is studied along two broad lines:
\begin{enumerate}
  \item \textbf{Classical geometry-driven methods:} explicit 2D/3D face models,
        landmarks, optical flow, or mesh-based warping
        \cite{thies2016face2face}. These favour interpretability and often run in
        real time when the representation is compact.
  \item \textbf{Neural generative methods:} models such as the First Order Motion
        Model (FOMM) \cite{siarohin2019first} or one-shot talking-head synthesis
        \cite{zakharov2020oneshot} learn highly nonlinear appearance--motion
        mappings and typically achieve more realistic texture and large head-pose
        variation, at the cost of heavier computation and weaker mathematical
        transparency.
\end{enumerate}

\paragraph{Landmark detectors.}
Reliable landmarks are a prerequisite for any geometry-based pipeline. Toolkits such
as OpenFace~\cite{baltrusaitis2016openface} and real-time mesh trackers in the
MediaPipe family~\cite{kartynnik2019attention} convert raw pixels into vectors in
$\mathbb{R}^{2k}$, which is exactly the interface to linear algebra used in the
present project.

\paragraph{Positioning of the chosen approach.}
The project adopts a \textbf{linear least-squares} map estimated from calibration
data, solved via \textbf{SVD}. This is closest in spirit to explicit geometry transfer
\cite{thies2016face2face} but simplified to a data-driven linear fit. A neural method
serves as an optional baseline for comparison on the same metrics.

\section{Chosen algorithm: outline, pros and cons}
\label{sec:algorithm}

\paragraph{Outline.}
\begin{enumerate}
  \item Detect $k$ landmarks on~$I_s$; form $x_s \in \mathbb{R}^{2k}$.
  \item \textbf{Calibration:} for selected frames~$t_i$, detect target landmarks
        $y_{t_i}$; build the system~\eqref{eq:Ma-b}.
  \item Compute $\hat{a} = \arg\min_a \|Ma-b\|_2^2$ using the SVD of~$M$.
  \item \textbf{Per frame:} detect~$y_t$, apply the fitted model to obtain
        predicted source landmarks~$\hat{y}_t$, warp~$I_s$ toward~$\hat{y}_t$,
        display the output.
\end{enumerate}

\paragraph{Pros.}
\begin{itemize}
  \item \textbf{Interpretability:} each step corresponds to standard linear-algebra
        objects (matrices, norms, subspaces, pseudoinverse).
  \item \textbf{Numerical stability:} SVD handles rank-deficient or
        ill-conditioned~$M$ more reliably than na\"ive normal equations.
  \item \textbf{Efficiency:} the dominant linear solve reuses a single factorisation
        computed during calibration; per-frame cost is dominated by detection and
        warping, not by the algebra.
\end{itemize}

\paragraph{Cons.}
\begin{itemize}
  \item \textbf{Linearity assumption:} real facial deformations and perspective
        effects are only approximately linear; large rotations and self-occlusion
        break the model.
  \item \textbf{Detector noise:} errors in landmark positions propagate directly
        into $M$ and~$b$, biasing~$\hat{a}$.
  \item \textbf{Limited expressiveness:} a single global linear map may underfit
        compared to deep models
        \cite{siarohin2019first,zakharov2020oneshot}.
\end{itemize}

\section{Theoretical background with emphasis on linear algebra}
\label{sec:theory}

\subsection{Landmarks as vectors and the least-squares problem}

All landmark configurations live in $V = \mathbb{R}^{2k}$. Stacking calibration
constraints produces $M \in \mathbb{R}^{m \times n}$ with $m \ge n$ and
$b \in \mathbb{R}^{m}$. The least-squares problem
\begin{equation}
  \min_{a \in \mathbb{R}^n} \|Ma - b\|_2^2
  \label{eq:ls}
\end{equation}
seeks the vector $Ma$ that is \emph{closest} to~$b$ in the Euclidean norm.

\begin{theorem}[Geometric characterisation]
Let $\mathrm{col}(M)$ denote the column space of~$M$. A vector $\hat{a}$ solves
\eqref{eq:ls} if and only if $M\hat{a}$ is the orthogonal projection of~$b$ onto
$\mathrm{col}(M)$; equivalently, the residual $r = b - M\hat{a}$ is orthogonal to
every column of~$M$:
\begin{equation}
  M^{\mathsf T}(b - M\hat{a}) = 0.
\end{equation}
\end{theorem}

\noindent
This yields the \emph{normal equations}
\begin{equation}
  M^{\mathsf T}M\,\hat{a} = M^{\mathsf T}b.
  \label{eq:normal}
\end{equation}
When $M$ has full column rank, $M^{\mathsf T}M$ is invertible and
$\hat{a} = (M^{\mathsf T}M)^{-1}M^{\mathsf T}b$. In floating-point arithmetic,
however, forming $M^{\mathsf T}M$ squares the condition number of~$M$, which is
undesirable when $M$ is ill-conditioned.

\subsection{Singular value decomposition and the pseudoinverse}

\begin{theorem}[SVD]
For any real $M \in \mathbb{R}^{m \times n}$ with $\mathrm{rank}(M)=r$, there exist
orthogonal matrices $U \in \mathbb{R}^{m \times m}$,
$V \in \mathbb{R}^{n \times n}$ and a diagonal
$\Sigma \in \mathbb{R}^{m \times n}$ with singular values
$\sigma_1 \ge \cdots \ge \sigma_r > 0$ such that
\begin{equation}
  M = U \Sigma V^{\mathsf T}.
\end{equation}
\end{theorem}

\noindent
Write $U = [U_r \; U_0]$ and $V = [V_r \; V_0]$ with
$U_r \in \mathbb{R}^{m \times r}$, $V_r \in \mathbb{R}^{n \times r}$ containing the
singular vectors corresponding to nonzero~$\sigma_i$. The
\emph{Moore--Penrose pseudoinverse} is then
\begin{equation}
  M^{+} = V_r \Sigma_r^{-1} U_r^{\mathsf T},
  \qquad
  \Sigma_r = \mathrm{diag}(\sigma_1,\ldots,\sigma_r).
\end{equation}
The least-squares solution of minimum Euclidean norm is
\begin{equation}
  \hat{a} = M^{+} b.
  \label{eq:pinv-ls}
\end{equation}
Thus SVD makes explicit the rank of~$M$, the subspaces $\mathrm{col}(M)$ and
$\ker(M)$, and provides a stable route to~$\hat{a}$ without explicitly
inverting~$M^{\mathsf T}M$.

\subsection{Conditioning and regularisation}

The sensitivity of~$\hat{a}$ to perturbations in~$b$ is governed by the singular
values: very small $\sigma_i$ amplify noise in the corresponding directions. The
\emph{condition number} $\kappa(M) = \sigma_1 / \sigma_r$ quantifies overall
sensitivity. In practice, \emph{truncated} or \emph{damped} pseudo\-inverses (e.g.\
Tikhonov regularisation) replace $\sigma_i^{-1}$ by
$\sigma_i / (\sigma_i^2 + \lambda^2)$, trading a small bias for reduced variance---a
standard linear-algebra viewpoint on ill-posed inverse problems.

\subsection{Connection to the application}

In the reenactment pipeline, $M$ encodes how calibration frames relate source and
target landmarks under the chosen linear parameterisation. \textbf{Low effective rank}
of~$M$ indicates that a few independent modes of motion dominate the calibration set;
\textbf{small singular values} flag directions where the data do not constrain~$a$,
which is expected when landmark noise is high or calibration diversity is insufficient.
SVD therefore serves both as a \emph{solver} and as a \emph{diagnostic} tool: it
reveals which degrees of freedom are well-determined and which require regularisation.

\section{Implementation pipeline}
\label{sec:pipeline}

\subsection{Data}

The input data consist of short webcam recordings (target) and a single high-resolution
portrait photograph (source). For each frame, a landmark detector extracts $k = 68$
facial key points. The calibration set is formed from approximately 50--100 frames
covering varied expressions (neutral, smiling, surprised, mouth open, eyebrows
raised). No external large-scale dataset is required; all data are collected by the
team and stored in the repository.

\subsection{Tools and libraries}

The implementation uses the \textbf{Python} programming language with the following
main libraries:
\begin{itemize}
  \item \textbf{MediaPipe} (or OpenFace) --- real-time facial landmark detection;
  \item \textbf{NumPy / SciPy} --- matrix operations and SVD computation
        (\texttt{numpy.linalg.svd}, \texttt{numpy.linalg.lstsq});
  \item \textbf{OpenCV} --- video capture, image warping (piecewise-affine or
        Delaunay triangulation), and display.
\end{itemize}

\subsection{Completed work}
\begin{itemize}
  \item Problem specification and mathematical formulation (landmark vectors,
        calibration system~\eqref{eq:Ma-b}, least-squares objective).
  \item Literature survey covering landmark-based, geometry-based, and neural
        reenactment approaches (see references).
  \item Design of evaluation metrics: mean Euclidean landmark error, normalised
        mean error (NME), cumulative error distribution (CED) curves with AUC,
        and image-quality measures (PSNR, SSIM).
  \item Pseudocode for pipeline~A (linear, SVD-based) and pipeline~B (neural
        baseline).
\end{itemize}

\subsection{Planned work}
\begin{itemize}
  \item Integrate the landmark detector and standardise the data layout for
        $P^{(S)}$ and $P^{(T)}$ matrices.
  \item Implement construction of $M$ and~$b$ from calibration frames and the
        SVD-based solve for~$\hat{a}$ (with optional truncation if small singular
        values are detected).
  \item Implement Delaunay-based piecewise-affine warping from source landmarks to
        predicted landmarks; profile end-to-end latency.
  \item Run a neural baseline (e.g.\ FOMM) on the same clips and report metrics
        side by side.
  \item Prepare numbered figures: calibration setup, example reenacted frames, CED
        curves, and a summary table of quantitative results.
\end{itemize}

\section{Preliminary results}
\label{sec:results}

At this stage the mathematical formulation is complete and verified on a toy example:
a synthetic set of $k=5$ landmarks with known affine transformation was used to
confirm that the SVD-based solver recovers the correct parameters~$\hat{a}$ up to
floating-point precision. On real webcam data, preliminary landmark extraction
produces visually plausible detections, but systematic quantitative evaluation
(error tables, CED curves, visual comparisons) is underway and will appear in the
final report.


\section{Conclusion}
\label{sec:conclusion}

This report formulates face reenactment as a least-squares landmark mapping
problem~\eqref{eq:ls} and shows that the SVD-based pseudoinverse~\eqref{eq:pinv-ls}
provides a numerically stable and mathematically transparent solution. A preliminary
toy-data test confirms that the solver correctly recovers known transformations.
The linear approach is expected to handle moderate expressions with low latency,
while large head rotations and self-occlusion are anticipated limitations where a
neural baseline may outperform the linear model. The remaining work focuses on
full integration of the landmark detector and image warping module, systematic
evaluation on real webcam data, and a side-by-side comparison with a neural method.

\section*{Link to code}

\noindent\textbf{GitHub repository:}
\texttt{\url{https://github.com/mykhailobrytan-dot/Project_Linear_Algebra}}

\begin{thebibliography}{9}

\bibitem{baltrusaitis2016openface}
T.~Baltru\v{s}aitis, P.~Robinson, and L.-P.~Morency,
\newblock OpenFace: an open source facial behavior analysis toolkit,
\newblock \emph{IEEE Winter Conference on Applications of Computer Vision (WACV)},
  2016.

\bibitem{kartynnik2019attention}
T.~Kartynnik, Y.~Ablavatski, I.~Grishchenko, and M.~Grundmann,
\newblock Real-time facial surface geometry from monocular video on mobile GPUs,
\newblock \emph{arXiv preprint arXiv:1907.06724}, 2019.

\bibitem{thies2016face2face}
J.~Thies, M.~Zollh\"ofer, M.~Stamminger, C.~Theobalt, and M.~Nie{\ss}ner,
\newblock Face2Face: real-time face capture and reenactment of RGB videos,
\newblock \emph{IEEE Conference on Computer Vision and Pattern Recognition (CVPR)},
  2016.

\bibitem{siarohin2019first}
A.~Siarohin, S.~Lathuili\`ere, S.~Tulyakov, E.~Ricci, and N.~Sebe,
\newblock First order motion model for image animation,
\newblock \emph{Advances in Neural Information Processing Systems (NeurIPS)}, 2019.

\bibitem{zakharov2020oneshot}
E.~Zakharov, A.~Ivakhnenko, A.~Shysheya, and V.~Lempitsky,
\newblock One-shot free-view neural talking-head synthesis for video conferencing,
\newblock \emph{IEEE Conference on Computer Vision and Pattern Recognition (CVPR)},
  2019.

\end{thebibliography}

\end{document}
